\newpage
\newcommand{\SpatialRows}{Homogeneous & Gaussian & 0.0130\% & 0.0000\% & 0.0046\% & 0.0025\% \\
Homogeneous & Near degeneracy & 0.0020\% & 0.0000\% & 0.0007\% & 0.0003\% \\
Inhomogeneous & Gaussian & 0.0001\% & 0.0000\% & 0.0000\% & 0.0000\% \\
Inhomogeneous & Near degeneracy & 0.0033\% & 0.0000\% & 0.0026\% & 0.0004\% \\}
\newcommand{\SpatialLocations}{Homogeneous & Gaussian & 0.0099887 & -20.7469 & 0.0000 & 0.00000 \\
Homogeneous & Near degeneracy & 0.0088803 & 7.7961 & 6.5381 & 0.00079 \\
Inhomogeneous & Gaussian & 0.0037238 & -8.9154 & 0.0000 & 0.00002 \\
Inhomogeneous & Near degeneracy & 0.0058621 & -841.2464 & 0.0000 & 0.00021 \\}
\newcommand{\SpatialLocationDrift}{0.00079}
\newcommand{\SpatialPeakDrift}{0.00459\%}
\newcommand{\SpatialQuadTighten}{$3.67\times10^{-12}\%$}
\newcommand{\SpatialBoundary}{0.00516\%}
\newcommand{\SpatialPassed}{4}

\section*{Stronger spatial reliability at $t=256$}
\textbf{Purpose.} Hold time, input, and $\alpha=1/2$ fixed, and vary the spatial
computation. This audit separates longitudinal enlargement from transverse
enlargement. It tests whether the four reported candidate maxima stabilize to
an exploratory relative tolerance of $0.1\%$.

The longitudinal period increases from 32,768 to 65,536 while the radial limit
stays at 2,048. We then search the additional annulus out to radius 4,096.
Finally, we double the longitudinal period again to 131,072, retaining the
larger radial limit. Every broad scan keeps $\Delta x=0.25$ and $\Delta r=2$.
The baseline is the saved final scan from the preceding experiment.

\begin{center}\small
\begin{tabular}{llrrrr}\toprule
Equation & Input & First $L$ change & Radius change & Second $L$ change & Integral check\\\midrule
\SpatialRows
\bottomrule\end{tabular}\end{center}
The three domain columns compare sampled peak magnitudes before and after the
stated expansion. The integral check compares the final Fourier-quadrature field
to an independently evaluated non-periodic integral at the refined candidate peak.
Percentages use the newer or independent result as the denominator. Values are rounded.

\textbf{Result.} \SpatialPassed\ of four cases pass the $0.1\%$ peak screen.
This includes the two longitudinal expansions, the transverse expansion, the
final locally optimized peak drift, and the independent integral check.
The largest change in locally optimized peak magnitude between the final two
periods is \SpatialPeakDrift.

\begin{center}\small
\begin{tabular}{llrrrr}\toprule
Equation & Input & Refined $R(256)$ & $x$ & $r$ & Location drift\\\midrule
\SpatialLocations
\bottomrule\end{tabular}\end{center}
Locations are continuous-coordinate candidate maxima from the independent
integral optimization. The drift is the Euclidean distance in $(x,r)$ between
local peak optimizations at the final two longitudinal periods. Its maximum is
\SpatialLocationDrift. Axial symmetry means a nonzero $r$ represents a ring of
equivalent points, not a unique three-dimensional location.

\textbf{Additional checks.} Local optimization starts from the broad search's
largest sample and refines both $x$ and $r$. Increasing the independent
integral's Gauss--Legendre order from 256 to 512 per panel changes the complex field by
at most \SpatialQuadTighten\ relative to the tighter result. Its integration
interval is split at the fractional multiplier cusp and into 64 Gaussian-frequency
panels. On each side, $k=-\beta\pm u^2$ removes the square-root cusp from the
integration variable. Fixed-order quadrature with an order-doubling check replaces
adaptive quadrature, which issued convergence warnings during local searches.
The largest sampled search-boundary amplitude is \SpatialBoundary\ of
the corresponding peak. Optimizer status and all intermediate measurements are
saved in the per-case JSON records.

\textbf{What we can now say.} At this time and for these four inputs, the peak
values have stronger numerical support from separate domain enlargements,
continuous local refinement, stable locations, and independent integration.
This improves confidence in the reported decay ratios. It does not prove the
maximum is global over all of $\mathbb R^3$, establish an asymptotic decay rate,
or extend reliability to other inputs, values of $\alpha$, or later times.

\textbf{Remaining spatial limitation.} The broad radial mesh is still spaced by
2; refinement is local to the identified maximum, not a global halving of every
cell. A missed narrow peak or a remote maximum is not excluded by these checks.
The report therefore claims numerical spatial stability of the identified peaks,
not a certified supremum.

\textbf{Reproduce.} Run \texttt{python3 spatial\_reliability.py}, then
\texttt{python3 report/make\_spatial\_tables.py}. Results, configurations, and
checkpoints are under \texttt{results/spatial\_reliability/}.
